\section{Detailed evaluation setup and statistical criteria}
\label{app:detailed-setup}

The Method defines the common answer scores, intervention objects, and restoration--deletion operations. This section fixes the model--interface pairs, evaluation examples, selection and confirmation roles, answer targets, reconstruction diagnostics, and statistical criteria used to apply those operations. The order matters: we first specify what is compared and which data may select a configuration, then define the diagnostics and uncertainty calculations applied without changing that configuration.

\subsection{Models and evaluation sets}

The protocol is instantiated in three frozen model--lens pairs. \textbf{\modelgemma{}} combines Gemma3 \citep{gemma3} with a per-layer transcoder (PLT), \textbf{\modelqwen{}} combines Qwen2.5-VL \citep{qwen25vl} with a PLT, and \textbf{\modelllava{}} combines LLaVA1.5 \citep{llava} with a cross-layer transcoder (CLT). These names describe our measurement setups, not official variants of the base models. Because the sparse interfaces are model-specific, cross-model tests align the image--question example, corruption, answer readout, intervention direction, and statistical orientation; they do not compare raw feature identities or raw score magnitudes. Native-channel experiments use coordinates from the frozen base models and do not use the transcoders.

We use three groups of evidence-region examples for different experimental roles. The first contains filtered OK-VQA validation examples \citep{okvqa} whose short answer can be determined from a localizable image region. Every example contains an image, a question, a fixed target answer, and a human-annotated evidence mask. The expanded evaluation contains 61 image--question examples. A separate 52-example evaluation is selected and annotated without using model results and serves as the independently selected confirmation for the behavior and hidden-position comparisons. Appendix~\ref{app:annotation-protocol} reports the sample accounting and annotation protocol for these compact-evidence sets.

The answer-position, same-object sparse, native-channel, and descriptive sparse-structure experiments use a second 300-image evidence set. Sixty images select layers, feature budgets, or channel fractions; the remaining 240 images evaluate those frozen choices. The same 240 image identifiers are used for all three models. They belong to 48 predefined groups of related entity--relation questions, which are resampled as clusters rather than treating every image--question row as independent. Because these images appeared in earlier project analyses, we call this a reused confirmation set rather than previously unseen held-out data. Its primary readout is the total log probability of every token in the fixed correct answer.

The third group supports the shared-subspace and full-hidden tests through three disjoint roles. Forty images fit each model's basis, 20 different images select the smallest candidate rank that meets the restoration and corrupted-state-replacement mean gates, and a separate 40-image set is reserved for replication of a selected rank. The tested ranks are $\mathcal{R}=\{1,2,4,8,16,32,64,128\}$, whose elements specify the number of retained basis directions. Earlier hidden-state screening fixed the layer and position group for each model before this subspace analysis: layer 1 at visual positions for Gemma, layer 14 at visual positions for Qwen, and layer 30 at the four answer-adjacent positions for LLaVA. The full-hidden reference is evaluated at the same interfaces. Any full-hidden interface whose mean effects meet the gate in both selection and replication is frozen for final confirmation on a separate pool of 61 candidate masks. Human review is completed without access to model predictions or intervention results, and only accepted masks enter the final confirmation.

A follow-up matched-interface analysis addresses the initial position asymmetry. On the same 60 accepted masks, all three models receive the same intervention at their zero-based penultimate language layer and four answer-adjacent positions: Gemma layer 32 of 34, Qwen layer 26 of 28, and LLaVA layer 30 of 32. The models use the same code, PyTorch runtime, answer-sequence scoring rule, and four paired gates. The shared-subspace test is also repeated at these matched relative locations on the same 40/20/40 fitting, selection, and replication image IDs. This relative-layer rule was chosen after the LLaVA result was known, so the follow-up is a matched diagnostic rather than a jointly prospective confirmation. Model architecture and tokenization remain different.

Answer targets are fixed before each paired comparison. Allowed short-answer aliases come from dataset answers and annotation metadata; once an alias is selected for an example, its tokenization is held constant across the paired corruptions and interventions. Empty targets, formatting failures, and cross-condition token mismatches are excluded from primary estimates and retained in the evaluation accounting. Sequence scores and correct-versus-wrong margins complement the first-token scores when an answer spans multiple tokens or has plausible alternatives.

\subsection{Sparse reconstruction diagnostics}

The same-object experiment first checks how accurately each sparse tool reconstructs the complete clean-minus-corrupted feed-forward-output change. We evaluate this reconstruction on the 60 development images at the four text positions immediately preceding the answer. Let $\mathcal{P}_{\mathrm{ans}}$ be this four-position set and $\mathcal{P}_{\mathrm{nz}}\subseteq\mathcal{P}_{\mathrm{ans}}$ the positions whose complete change is nonzero. The relative squared reconstruction error and mean direction similarity are
\begin{align}
R_{\mathrm{recon}}(x,c)
&=\frac{\sum_{i\in\mathcal{P}_{\mathrm{ans}}}\|\vec{r}_{\mathrm{err},i}^{m}(x,c;\theta,\phi)\|_2^2}
        {\sum_{i\in\mathcal{P}_{\mathrm{ans}}}\|\Delta\vec{u}_{i}^{m}(x,c;\theta)\|_2^2},
\label{eq:relative-reconstruction-error}\\
C_{\mathrm{dir}}(x,c)
&=\frac{1}{|\mathcal{P}_{\mathrm{nz}}|}
  \sum_{i\in\mathcal{P}_{\mathrm{nz}}}
  \frac{\left\langle\widehat{\Delta\vec{u}}_{K,i}^{\ell\rightarrow m},
  \Delta\vec{u}_{i}^{m}\right\rangle}
  {\left\|\widehat{\Delta\vec{u}}_{K,i}^{\ell\rightarrow m}\right\|_2
   \left\|\Delta\vec{u}_{i}^{m}\right\|_2}.
\label{eq:reconstruction-direction-similarity}
\end{align}
The cosine uses the same $(x,c;\theta,\phi)$ arguments as the reconstruction in Equation~\ref{eq:sparse-reconstruction}; these repeated arguments are omitted inside Equation~\ref{eq:reconstruction-direction-similarity} for readability. A zero-norm reconstruction contributes zero similarity, while an image with no complete change at any of the four positions is excluded from the image-level mean. We first compute each scalar within an image and then average across valid development images. A value $R_{\mathrm{recon}}>1$ means that the squared reconstruction error exceeds the squared size of the complete change, while larger $C_{\mathrm{dir}}$ means closer directional agreement.

\subsection{Statistical inference and decision criteria}

The image--question example is the inferential unit. Multiple prompt, mask, feature, or intervention rows belonging to the same example are first averaged within that example; 95\% bootstrap confidence intervals then resample image--question examples rather than raw rows. We report raw row counts only as computational coverage. The 61- and 52-example panel intervals are unadjusted across cells: an interval above zero describes a positive individual contrast, not a family-wise decision. A positive estimate whose interval crosses zero is reported as directional.

For the 240-image experiments, resampling and sign-flip tests operate on the 48 entity--relation clusters. Confidence intervals use 10,000 paired cluster-bootstrap draws and tests use 100,000 sign flips. The answer-position family corrects 12 comparisons together. The sparse-interface family jointly corrects 72 two-sided tests: 36 component effects (three models, three position groups, two components, two directions), 12 random- or wrong-position contrasts (three models, two controls, two directions), and 24 cross-model structural contrasts (three model pairs, two position groups, four summaries). The three component-effect groups are visual positions, answer-adjacent text positions, and their combined set; Figure~\ref{fig:sparse-component-effects} plots only the combined group. This single Holm correction spans effect, control, and structural questions; the 16-, 32-, and 64-feature budget curves are descriptive and are not part of the 72 tests. The native-channel family corrects 66 tests: three models times nine contrasts under positive-only selection and 13 under positive-and-negative selection. A contrast passes only when it meets its predeclared direction or two-sided condition after the corresponding joint correction.

The shared-subspace test uses the complete-answer sequence score and image--question examples as the inferential unit. For a candidate rank, both the mean restoration and mean corrupted-state-replacement effect must fall between 80\% and 120\% of the mean clean-minus-masked sequence-score change on the 20-image selection set. This reference is the clean-to-evidence-mask damage in Equation~\ref{eq:corruption-damage}, not the evidence-versus-control contrast in Equation~\ref{eq:evidence-control-contrast}. Only the smallest passing rank advances to the 40-image replication set, where four paired sample-bootstrap contrasts test the lower and upper bounds for the two directions. A rank that fails selection does not receive strict confirmation. The full-hidden reference uses the same mean gate to nominate one frozen interface. Its independent confirmation requires the 95\% sample-bootstrap lower bounds of all four contrasts to be nonnegative; the mask review, interface, and decision rule are frozen before model evaluation. The matched follow-up applies these four bounds to each model's 60 rows without using its results to rescan layers or positions.

Cross-model conclusions compare score-change direction, positive-sample fraction, within-model retention, and gate status. Raw score changes are not ranked across models because architectures, vocabularies, layers, and intervention coordinates differ. Appendix~\ref{app:sparse-lens-details} records lens checkpoints, layers, and feature counts; Appendix~\ref{app:shared-subspace-details} reports the complete rank sweep and confirmation accounting; and Appendix~\ref{app:detailed-method} defines the intervention procedure.
